% Rozdziały dokumentujące pracę własną studenta: opisujące ideę, sposób lub metodę 
% rozwiązania postawionego problemu oraz rozdziały opisujące techniczną stronę rozwiązania 
% --- dokumentacja techniczna, przeprowadzone testy, badania i uzyskane wyniki. 

% Praca musi zawierać elementy pracy własnej autora adekwatne do jego wiedzy praktycznej uzyskanej w
% okresie studiów. Za pracę własną autora można uznać np.: stworzenie aplikacji informatycznej lub jej
% fragmentu, zaproponowanie algorytmu rozwiązania problemu szczegółowego, przedstawienie projektu 
% np.~systemu informatycznego lub sieci komputerowej, analizę i ocenę nowych technologii lub rozwiązań
% informatycznych wykorzystywanych w przedsiębiorstwach, itp. 

% Autor powinien zadbać o właściwą dokumentację pracy własnej obejmującą specyfikację założeń i 
% sposób realizacji poszczególnych zadań
% wraz z ich oceną i opisem napotkanych problemów. W przypadku prac o charakterze 
% projektowo-implementacyjnym, ta część pracy jest zastępowana dokumentacją techniczną i użytkową systemu. 

% W pracy \textbf{nie należy zamieszczać całego kodu źródłowego} opracowanych programów. Kod źródłowy napisanych
% programów, wszelkie oprogramowanie wytworzone i wykorzystane w pracy, wyniki przeprowadzonych
% eksperymentów powinny być umieszczone np. na płycie CD, stanowiącej dodatek do pracy.

% \section*{Styl tekstu}

% Należy\footnote{Uwagi o stylu pochodzą częściowo ze stron prof. Macieja Drozdowskiego~\cite{Drozdowski2006}.} 
% stosować formę bezosobową, tj.~\emph{w pracy rozważono ......, 
% w ramach pracy zaprojektowano ....}, a nie: \emph{w pracy rozważyłem, w ramach pracy zaprojektowałem}. 
% Odwołania do wcześniejszych fragmentów tekstu powinny mieć następującą postać: ,,Jak wspomniano wcześniej, ....'', 
% ,,Jak wykazano powyżej ....''. Należy unikać długich zdań. 

% Niedopuszczalne są zwroty używane w języku potocznym. W pracy należy używać terminologii informatycznej, która ma 
% sprecyzowaną treść i znaczenie. 

% Niedopuszczalne jest pisanie pracy metodą \emph{cut\&paste}, bo jest to plagiat i dowód intelektualnej indolencji autora.
% Dane zagadnienie należy opisać własnymi słowami. Zawsze trzeba powołać się na zewnętrzne źródła. 

% ////////////////////////////////////////////////////////////////////////////////////////
% 1. Interfejsy i testy bazowe
% • Implementacja bazowego projektu neoTRNG na układzie FPGA.
% • Uruchomienie komunikacji:
% • interfejs UART
% • moduł pMOD SD Card (szybki zapis surowych bitów bezpośrednio do
% sektorów karty, bez systemu plików, przez interfejs SPI).
% • Przygotowanie środowiska analitycznego na komputerze PC: instalacja i
% konfiguracja pakietu ent, zestawu testów statystycznych NIST SP 800-22 oraz
% narzędzia do estymacji min-entropii NIST SP 800-90B. (bez/z de-biassing, bez/z
% CRC (z obecnym de-biasing))

\chapter{Testy bazowe}

\noindent W niniejszym rozdziale przeanalizowano wpływ zastosowania korekcji źródła oraz 
przetwarzania końcowego na poziom generowanej entropii w domyślnej implementacji 
neoTRNG. Rozważane elementy architektury zostały przedstawione na rysunku 
\ref{fig:architektura_neotrng} jako moduły „de-biasing” oraz „sampling”. W 
implementacji bazowej zastosowano korektor von Neumanna oraz mechanizm kompresji 
entropii oparty na CRC. W ramach przeprowadzonych testów porównano cztery 
warianty konfiguracji generatora:

\begin{enumerate}
    \item moduł z korektorem von Neumanna oraz CRC (domyślny),
    \item moduł z korektorem von Neumanna bez CRC,
    \item moduł bez korektora z aktywnym CRC,
    \item moduł bez korektora i bez CRC.
\end{enumerate}

Zbadane zostały kolejno moduły z:

\begin{enumerate}
    \item trzema pętlami oscylatorów po 5, 7 oraz 9 inwerterów,
    \item trzema pętlami oscylatorów po 3, 5 oraz 7 inwerterów,
    \item dwiema pętlami oscylatorów po 3 i 5 inwerterów.
\end{enumerate}

\section{Modyfikacja}

\noindent Testy przeprowadzono na zmodyfikowanym module neoTRNG. Moduł rozszerzono o dwa parametry generyczne
 \textit{ENABLE\_VON\_NEUMANN} oraz \textit{ENABLE\_CRC} w celu łatwej konfiguracji. Zmodyfikowany plik źródłowy znajduje się w załączniku
  A neoTRNG\_mod\_1.vhd.

Każdy z przypadków testowych uruchomiono na sprzęcie, a uzyskane wyniki zapisano w odpowiednich plikach.


\section{ENT}

\begin{table}[H]
\centering
\caption{ENT: Entropy bits per byte.}
\label{tab:ch3_ent}
\begin{tabular}{|l|c|c|c|c|}
\hline
\textbf{Konfiguracja} & \textbf{3RO 5-7-9} & \textbf{3RO 3-5-7} & \textbf{2RO 3-5} \\ \hline
\textbf{true/true}      & 7.999998 & 7.999998 & 7.999997 \\ \hline
\textbf{true/false}     & 7.928600 & 7.856006 & 7.368478 \\ \hline
\textbf{false/true}     & 7.999988 & 7.999995 & 7.999742 \\ \hline
\textbf{false/false}    & 7.404643 & 7.921076 & 7.617650 \\ \hline
\end{tabular}
\end{table}

\section{NIST SP 800-22}

\begin{flushleft}
\small \textbf{Note:} The minimum pass rate for each statistical test with the exception of the
random excursion (variant) test is approximately = 96 for a
sample size = 100 binary sequences\cite{NIST-SP-800-22}. \\
\small \textbf{Note:} The minimum pass rate for each statistical test with the exception of the
random excursion (variant) test is approximately = 96 for a
sample size = 100 binary sequences\cite{NIST-SP-800-22}. \\
\small \textbf{Legend:} \textbf{sp} - subtests passed (zaliczone podtesty). \\
\small \textbf{Legend:} \textbf{*} - test niezaliczony. \\
\small \textbf{Legend:} \textbf{\textit{ENABLE\_VON\_NEUMANN}/ \textit{ENABLE\_CRC}} - wybrana konfiguracja.
\end{flushleft}

\begin{table}[H]
\centering
\caption{3 RO : 5 - 7 - 9}
\label{tab:nist_22_3RO_21I}
\begin{tabular}{|l|c|l|c|l|c|l|c|l|}
\hline
\textbf{Statistical Test} & \multicolumn{2}{c|}{\textbf{true/true}} & \multicolumn{2}{c|}{\textbf{true/false}} & \multicolumn{2}{c|}{\textbf{false/true}} & \multicolumn{2}{c|}{\textbf{false/false}} \\ \hline
Frequency               &  98/100 &     & 100/100 &     &  97/100 &     & 100/100 &     \\ \hline
BlockFrequency          &  99/100 &     & 100/100 &     &  98/100 &     & 100/100 &     \\ \hline
CumulativeSums          &  98/100 &     & 100/100 &     &  97/100 &     &  99/100 &     \\ \hline
CumulativeSums          &  98/100 &     & 100/100 &     &  97/100 &     & 100/100 &     \\ \hline
Runs                    & 100/100 &     &   0/100 & *   &  97/100 &     &   0/100 & *   \\ \hline
LongestRun              & 100/100 &     &   0/100 & *   &  99/100 &     &   0/100 & *   \\ \hline
Rank                    &  99/100 &     &  99/100 &     &  99/100 &     &  99/100 &     \\ \hline
FFT                     & 100/100 &     &   2/100 & *   &  98/100 &     &   0/100 & *   \\ \hline
NonOverlappingTemplate  & 147/148 & sp  &   8/148 & sp  & 148/148 & sp  &   0/148 & sp  \\ \hline
OverlappingTemplate     & 100/100 &     &   0/100 & *   &  99/100 &     &   0/100 & *   \\ \hline
Universal               &  97/100 &     &   0/100 & *   & 100/100 &     &   0/100 & *   \\ \hline
ApproximateEntropy      & 100/100 &     &   0/100 & *   & 100/100 &     &   0/100 & *   \\ \hline
RandomExcursions        &     7/8 & sp  &     1/8 & sp  &     8/8 & sp  &     2/8 & sp  \\ \hline
RandomExcursionsVariant &   18/18 & sp  &   18/18 & sp  &   18/18 & sp  &   18/18 & sp  \\ \hline
Serial                  &  99/100 &     &   0/100 & *   &  99/100 &     &   0/100 & *   \\ \hline
Serial                  &  98/100 &     &  95/100 & *   &  99/100 &     &   0/100 & *   \\ \hline
LinearComplexity        &  99/100 &     &  99/100 &     & 100/100 &     & 100/100 &     \\ \hline
\end{tabular}
\end{table}

\begin{table}[H]
\centering
\caption{3 RO : 3 - 5 - 7}
\label{tab:nist_22_3RO_15I}
\begin{tabular}{|l|c|l|c|l|c|l|c|l|}
\hline
\textbf{Statistical Test} & \multicolumn{2}{c|}{\textbf{true/true}} & \multicolumn{2}{c|}{\textbf{true/false}} & \multicolumn{2}{c|}{\textbf{false/true}} & \multicolumn{2}{c|}{\textbf{false/false}} \\ \hline
Frequency               &  98/100 &     &  96/100 &     &  98/100 &     &  94/100 & *   \\ \hline
BlockFrequency          & 100/100 &     &   0/100 & *   &  97/100 &     &   0/100 & *   \\ \hline
CumulativeSums          &  98/100 &     &  93/100 & *   &  99/100 &     &  93/100 & *   \\ \hline
CumulativeSums          &  99/100 &     &  93/100 & *   &  98/100 &     &  95/100 & *   \\ \hline
Runs                    & 100/100 &     &   0/100 & *   &  98/100 &     &  84/100 & *   \\ \hline
LongestRun              &  99/100 &     &   0/100 & *   &  99/100 &     &  77/100 & *   \\ \hline
Rank                    &  99/100 &     &  98/100 &     & 100/100 &     &  97/100 &     \\ \hline
FFT                     &  98/100 &     &   0/100 & *   &  99/100 &     &   7/100 & *   \\ \hline
NonOverlappingTemplate  & 147/148 & sp  &   0/148 & sp  & 148/148 & sp  &  11/148 & sp  \\ \hline
OverlappingTemplate     & 100/100 &     &   0/100 & *   &  99/100 &     &   1/100 & *   \\ \hline
Universal               &  98/100 &     &   0/100 & *   &  96/100 &     &   0/100 & *   \\ \hline
ApproximateEntropy      & 100/100 &     &   0/100 & *   &  99/100 &     &   0/100 & *   \\ \hline
RandomExcursions        &     8/8 & sp  &     0/8 & sp  &     8/8 & sp  &     4/8 & sp  \\ \hline
RandomExcursionsVariant &   18/18 & sp  &   18/18 & sp  &   18/18 & sp  &   18/18 & sp  \\ \hline
Serial                  &  99/100 &     &   0/100 & *   &  98/100 &     &   0/100 & *   \\ \hline
Serial                  &  99/100 &     &  20/100 & *   &  99/100 &     &  94/100 & *   \\ \hline
LinearComplexity        &  97/100 &     &  98/100 &     & 100/100 &     &  99/100 &     \\ \hline
\end{tabular}
\end{table}

\begin{table}[H]
\centering
\caption{2 RO : 3 - 5}
\label{tab:nist_22_2RO_8I}
\begin{tabular}{|l|c|l|c|l|c|l|c|l|}
\hline
\textbf{Statistical Test} & \multicolumn{2}{c|}{\textbf{true/true}} & \multicolumn{2}{c|}{\textbf{true/false}} & \multicolumn{2}{c|}{\textbf{false/true}} & \multicolumn{2}{c|}{\textbf{false/false}} \\ \hline
Frequency               & 100/100 &     & 100/100 &     &  98/100 &     &  14/100 & *   \\ \hline
BlockFrequency          & 100/100 &     & 100/100 &     &  98/100 &     &   0/100 & *   \\ \hline
CumulativeSums          & 100/100 &     & 100/100 &     &  98/100 &     &  14/100 & *   \\ \hline
CumulativeSums          & 100/100 &     & 100/100 &     &  98/100 &     &  15/100 & *   \\ \hline
Runs                    & 100/100 &     &   0/100 & *   &  99/100 &     &   0/100 & *   \\ \hline
LongestRun              &  98/100 &     &   0/100 & *   &  99/100 &     &   0/100 & *   \\ \hline
Rank                    &  99/100 &     &  99/100 &     &  99/100 &     &  99/100 &     \\ \hline
FFT                     & 100/100 &     &   0/100 & *   &  98/100 &     &   0/100 & *   \\ \hline
NonOverlappingTemplate  & 148/148 & sp  &   8/148 & sp  & 148/148 & sp  &   4/148 & sp  \\ \hline
OverlappingTemplate     &  99/100 &     &   0/100 & *   & 100/100 &     &   0/100 & *   \\ \hline
Universal               &  98/100 &     &   0/100 & *   & 100/100 &     &   0/100 & *   \\ \hline
ApproximateEntropy      &  98/100 &     &   0/100 & *   &  97/100 &     &   0/100 & *   \\ \hline
RandomExcursions        &     8/8 & sp  &     0/8 & sp  &     8/8 & sp  &     0/8 & sp  \\ \hline
RandomExcursionsVariant &   18/18 & sp  &   18/18 & sp  &   18/18 & sp  &   18/18 & sp  \\ \hline
Serial                  &  99/100 &     &   0/100 & *   &  98/100 &     &   0/100 & *   \\ \hline
Serial                  & 100/100 &     &   0/100 & *   &  99/100 &     &   0/100 & *   \\ \hline
LinearComplexity        &  98/100 &     & 100/100 &     & 100/100 &     & 100/100 &     \\ \hline
\end{tabular}
\end{table}



\section{NIST SP 800-90B}



\begin{table}[H]
\centering
\caption{NIST SP 800-90: Non-IID entropy assessment}
\label{tab:ch3_NIST_SP_800_90}
\begin{tabular}{|l|c|c|c|c|}
\hline
\textbf{Konfiguracja} & \textbf{3RO 5-7-9} & \textbf{3RO 3-5-7} & \textbf{2RO 3-5} \\ \hline

\textbf{true/true} &  
\shortstack{$H_{orig}=7.918$ \\ $H_{bit}=0.941$} & 
\shortstack{$H_{orig}=$ \\ $H_{bit}=$} & 
\shortstack{$H_{orig}=$ \\ $H_{bit}=$} \\ \hline

\textbf{true/false} &
\shortstack{$H_{orig}=6.653213$ \\ $H_{bit}=0.754510$} &
\shortstack{$H_{orig}=$ \\ $H_{bit}=$} &
\shortstack{$H_{orig}=$ \\ $H_{bit}=$} \\ \hline

\textbf{false/true} &
\shortstack{$H_{orig}=7.915044$ \\ $H_{bit}=0.941182$} &
\shortstack{$H_{orig}=$ \\ $H_{bit}=$} &
\shortstack{$H_{orig}=$ \\ $H_{bit}=$} \\ \hline

\textbf{false/false} &
\shortstack{$H_{orig}=$ \\ $H_{bit}=$} &
\shortstack{$H_{orig}=$ \\ $H_{bit}=$} &
\shortstack{$H_{orig}=$ \\ $H_{bit}=$} \\ \hline
\end{tabular}
\end{table}













